% Nouveau EGA V (cote 25): archive cote 25, site page 149 (PDF page 150).
% typescript. Draft compared: none.
% AI-transcribed by Claude Opus 5.5 (claude-opus-5-5, adaptive thinking, effort high), Message Batches API, batch ega25b (msgbatch_01A7n5YuXubUPNK1zsepi8RG), 2026-11-13; not yet checked by eye.
\scanpage{25/149}
\mspage{10}
% NOTE: a long handwritten diagonal stroke crosses the whole text, from ``Arrivé'' at the top down to the foot of the page. It may cancel the page or mark it as used. The text under it is transcribed as live text.

Arrivé à ce point \unsure{admirable|admitable} de généralité, je m'aperçois malheureuse-%
ment que ces développements sont encore trop particuliers, et qu'au lieu
d'un sous-préschéma défini par une seule section de $\underline{O}_X$ , il y a lieu de
considérer celui défini par m sections, ou ce qui\hand{/}revient au même, par un
homomorphisme $\underline{E} \longrightarrow \underline{O}_X$ , où $\underline{E}$ est un Module loc libre de rang m. Je te
laisse le soin de décider s'il y a lieu dès le début de commencer avec le
cas général, suivant le principe qu'il faut procéder du général au parti-%
culier.

Introduisons d'abord une terminologie pour des anneaux A à
augmentation $A \to A/I$ (qui sera utilisée également dans le contexte
faisceautisé). L'augmentation est dite \uline{régulière} (ou devons nous dire
quasi-régulière, avec les définitions adoptées au par.19 ?) si \inserted{a)} $I/I^2$ est
un A/I -Module projectif de type fini, et \inserted{b)} l'homomorphisme
\[ \theta\colon \mathrm{Sym}\,(I/I^2) \longrightarrow \mathrm{gr}_I(A) \]
est un isomorphisme. Nous dirons que l'augmentation est \inserted{de type} ``\uline{quadratique}
\uline{ordinaire}'' lorsque a) ci-dessus est vérifié, ainsi que b') l'homomor-%
phisme canonique précédent à un noyau engendré par $\theta^2$ , lequel noyau $\underline{K}$
est un \inserted{sous-}Module inversible ``non dégénéré'' de $\mathrm{Sym}^2(I/I^2)$ . Par ``non dégé%
néré'' nous entendons que l'homomorphisme correspondant
\[ \underline{K} \otimes (\underline{I}/\underline{I}^2)^{\hand{\vee}} \longrightarrow \underline{I}/\underline{I}^2 \]
(déduit en utilisant l'homomorphisme canonique
\[ \mathrm{Sym}^2(I/I^2) \longrightarrow \mathrm{Hom}\ (\ (I/I^2)^{\hand{\vee}}\ ,\ I/I^2\ )\quad ) \]
est bijectif, ou ce qui revient au même surjectif (car il s'agit de
Modules projectifs \inserted{sur A/I} de même rang en chaque point), ou encore : prenant
localement une base \inserted{$\phi$} de $\underline{K}$ , \xout{l'homomorphisme correspondant} la forme
quadratique correspondante sur $(I/I^2)^{\hand{\vee}}$ est non dégénérée. Il suffit

\mspage{149}
